\documentclass[11pt]{article}
\usepackage{mathblatt}
% gesetzt von werkzeuge/setzer.py; Lösungen nach bau/bauregeln.md „Lösungen“.
\newcommand{\kennung}{FEU}
\begin{document}
\pfheftstil
\fancyfoot[C]{{\footnotesize\color{mbgrau}\kennung\hfill\thepage}}
\pfheftkopf{Quader und Würfel}{Klasse 6 \textperiodcentered{} Lösungen}

\begin{pfloesung}{}
\lz{1b)}{$8$ Würfel, $3$ Schichten, $V = 24$\,cm³}{$4 \cdot 2 = 8$ Würfel je Schicht; $3$ Schichten; $8 \cdot 3 = 24$\,cm³}{}
\end{pfloesung}
\begin{pfloesung}{}
\lz{2.}{$V = 84$\,cm³}{$V = 7 \cdot 4 \cdot 3 = 84$\,cm³}{}
\end{pfloesung}
\begin{pfloesung}{}
\lz{3.}{$V = 216$\,cm³}{$V = 6 \cdot 6 \cdot 6 = 216$\,cm³}{}
\end{pfloesung}
\begin{pfloesung}{}
\lz{4.}{In den Würfel, $53$\,dm³ mehr}{Kiste: $6 \cdot 4 \cdot 3 = 72$\,dm³; Würfel: $5 \cdot 5 \cdot 5 = 125$\,dm³; $125 - 72 = 53$\,dm³}{}
\end{pfloesung}
\begin{pfloesung}{}
\lz{5b)}{$O = 112$\,cm²}{$8 \cdot 4 = 32$; $8 \cdot 2 = 16$; $4 \cdot 2 = 8$; $O = 2 \cdot 32 + 2 \cdot 16 + 2 \cdot 8 = 112$\,cm²}{}
\end{pfloesung}
\begin{pfloesung}{}
\lz{6.}{$O = 600$\,cm²}{$10 \cdot 10 = 100$\,cm²; $O = 6 \cdot 100 = 600$\,cm²}{}
\end{pfloesung}
\begin{pfloesung}{}
\lz{7.}{$56$\,dm²}{Boden $5 \cdot 4 = 20$; vorn und hinten $2 \cdot 5 \cdot 2 = 20$; Seiten $2 \cdot 4 \cdot 2 = 16$; $20 + 20 + 16 = 56$\,dm²}{}
\end{pfloesung}
\begin{pfloesung}{}
\lz{8.}{Nein, er braucht $9\,000$\,cm²}{Nein; Boden $60 \cdot 30 = 1\,800$; vorn und hinten $2 \cdot 60 \cdot 40 = 4\,800$; Seiten $2 \cdot 30 \cdot 40 = 2\,400$; $1\,800 + 4\,800 + 2\,400 = 9\,000$\,cm²; es fehlen $9\,000 - 8\,000 = 1\,000$\,cm²}{}
\end{pfloesung}
\begin{pfloesung}{}
\lz{9b)}{$5$ Liter; $3\,000$\,cm³; $2{,}5$ Liter; $4\,000$ Liter}{$5\,000 : 1\,000 = 5$ Liter; $3 \cdot 1\,000 = 3\,000$\,cm³; $2\,500 : 1\,000 = 2{,}5$ Liter; $4 \cdot 1\,000 = 4\,000$ Liter}{}
\end{pfloesung}
\begin{pfloesung}{}
\lz{10.}{$6$ Liter}{$V = 30 \cdot 20 \cdot 10 = 6\,000$\,cm³ $= 6$ Liter}{}
\end{pfloesung}
\begin{pfloesung}{}
\lz{11.}{$6\,000$ Liter}{$V = 3 \cdot 2 \cdot 1 = 6$\,m³; $6 \cdot 1\,000 = 6\,000$ Liter}{}
\end{pfloesung}
\begin{pfloesung}{}
\lz{12.}{Nein, weniger: $63$ Liter}{$V = 60 \cdot 30 \cdot 35 = 63\,000$\,cm³ $= 63$ Liter}{}
\end{pfloesung}
\begin{pfloesung}{}
\lz{13b)}{$3$\,cm hoch}{Grundfläche $5 \cdot 4 = 20$\,cm²; $60 : 20 = 3$\,cm}{}
\end{pfloesung}
\begin{pfloesung}{}
\lz{14.}{$20$\,cm hoch}{$12$ Liter $= 12\,000$\,cm³; $30 \cdot 20 = 600$\,cm²; $12\,000 : 600 = 20$\,cm}{}
\end{pfloesung}
\begin{pfloesung}{}
\lz{15.}{$20$\,cm hoch}{$30$ Liter $= 30\,000$\,cm³; $60 \cdot 25 = 1\,500$\,cm²; $30\,000 : 1\,500 = 20$\,cm}{}
\end{pfloesung}
\begin{pfloesung}{}
\lz{16.}{Ja, $10$\,cm bleiben frei}{Ja; $98$ Liter $= 98\,000$\,cm³; $70 \cdot 40 = 2\,800$\,cm²; $98\,000 : 2\,800 = 35$\,cm; frei $45 - 35 = 10$\,cm}{}
\end{pfloesung}
\begin{pfloesung}{}
\lz{17.}{$64$\,cm³}{$4 \cdot 4 \cdot 4 = 64$ (nicht $4 \cdot 4$, nicht die Oberfläche $96$)}{}
\end{pfloesung}
\begin{pfloesung}{}
\lz{18.}{$6{,}5$ Liter}{$6\,500 : 1\,000 = 6{,}5$ Liter}{}
\end{pfloesung}
\begin{pfloesung}{}
\lz{19.}{$V = 90$\,cm³, $O = 126$\,cm²}{$V = 6 \cdot 5 \cdot 3 = 90$\,cm³; $O = 2 \cdot 30 + 2 \cdot 18 + 2 \cdot 15 = 126$\,cm²}{}
\end{pfloesung}
\begin{pfloesung}{}
\lz{20.}{$10$\,cm hoch}{$2$ Liter $= 2\,000$\,cm³; $20 \cdot 10 = 200$\,cm²; $2\,000 : 200 = 10$\,cm}{}
\end{pfloesung}
\begin{pfloesung}{}
\lz{21.}{Ja, man braucht $96$ Liter}{Ja; Wasserhöhe $50 - 10 = 40$\,cm; $80 \cdot 30 \cdot 40 = 96\,000$\,cm³ $= 96$ Liter; $100 - 96 = 4$ Liter bleiben übrig}{}
\end{pfloesung}
\end{document}
